\renewcommand{\theposttitle}{Authenticity}
\renewcommand{\thepostdate}{2023-05-30}
\vspace*{\fill}
\begin{center}
\includegraphics[width=\textwidth, height=0.55\textheight, keepaspectratio]{C:/Users/James/Documents/GitHub/VisualIdeasBook/data/images/f4bccdc24537d9b0b6cf175e8bc9325e.png}
\end{center}
\vspace{2mm}
\begin{center}
\begin{minipage}{0.92\textwidth}
\small
\textbf{It is the cracked pot that lets the light in.}\\
-Leonard Cohen

Not a direct quote, but a paraphrased one.

This image drew up some challenging comments on Instagram. A few people remarked how they were betrayed by revealing too much to the people they trusted.

This reminded me how I’ve been burned too, more times than I can count. In fact, the one thing I remember from my dad growing up is how many times his friends swindled him out of his hard-earned money because he revealed too many scars, which created real fears in me around trust.

But I’ve also been lucky to meet people who restored some of that trust. \\
• An art teacher who encouraged me out of a sea of 4,000 other low-income, public school students to keep creating. \\
• An elderly Japanese man in Hiroshima who saw me lost in the train station and spent 2 hours to help me get to where I was going. \\
• A friend I just met in Thailand inviting me over to share in Christmas dinner and presents so I wouldn’t spend the holidays alone.

Some scars heal and some scars fester. All I know is I saw what the festering did to my dad and that’s something I’d like to avoid.

I’m curious, how do you navigate revealing who you are on social media?
\end{minipage}
\end{center}
\vspace*{\fill}